\hspace{3mm}

\centerline{\bf \Huge Математическая драка}

\problem{1}
Маша заменила в примере на умножение двузначных чисел цифры буквами: одинаковые – одинаковыми, разные – разными. У неё получилось АБ$\cdot$ВГ = БББ Каким мог быть исходный пример? Найдите все возможные варианты.

\problem{2}
Оля задумала четыре целых числа, а затем нашла все их попарные суммы. Пять из них оказались равны 70, 110, 120, 180 и 230. Чему равна шестая сумма?

\problem{3}
Все четырехзначные числа, цифры которых различны и стоят в порядке возрастания, выписали друг за другом — снова в порядке возрастания. Какое число стоит на 97–м месте?

\problem{4}
На карточке у каждого из 11 школьников написано какое–то натуральное число. Известно, что все эти числа различны, а их сумма равна 71. Чему может быть равно среднее по величине из написанных чисел?

\problem{5}
Сколько существует пятизначных чисел, десятичная запись которых начинается с 1 и содержит ровно две одинаковые цифры?

\problem{6}
Пусть A — двузначное число, не кратное 10. B — трехзначное число. Известно, что A\% от B равны 400. Найдите, чему могут быть равны А и B.

\problem{7}
Грузчики Коля и Петя носят ящики. Переноска маленького ящика занимает у Пети 1 минуту, а у Коли 3 минуты. Зато большой ящик Коля переносит за 5 минут, а Петя — за 6. Всего им нужно перенести 10 больших и 10 маленьких ящиков. За какое наименьшее время они могут это сделать?

\problem{8}
Найдите 2009–й член последовательности: $1, 2, 1, 2, 3, 2, 1, 2, 3, 4, 3, 2, 1, 2, 3, 4, 5, 4, 3, 2, 1\ldots$
